%  01_introduction.tex
\section{Introduction}
\label{sec:intro}

Fruit bats play pivotal roles in tropical and subtropical ecosystems as pollinators and
seed dispersers, yet the individual-level behaviour of colony members remains poorly
understood.  Long-term studies of social dynamics, foraging ranges, and roosting fidelity
require reliable identification of individual animals across repeated observations---a
task traditionally performed by physical tagging or manual expert annotation, both of
which are costly, disruptive to the animals, and difficult to scale.

Automated re-identification from facial imagery offers a non-invasive alternative.
The bat face is sufficiently individuated through pelage patterns, vein geometry, and
facial proportions to support individual recognition~\cite{ravela2004visual}.
Advances in deep face recognition for humans~\cite{parkhi2015deep, schroff2015facenet,
deng2019arcface, kim2022adaface} provide powerful transfer-learning priors, yet animal
datasets pose three challenges absent in standard human benchmarks:
(i)~\emph{small identity count}---studies typically monitor tens, not millions, of individuals;
(ii)~\emph{open-set inference}---a deployed system must handle bats not seen during training;
and (iii)~\emph{environmental variation}---imaging conditions change across sessions,
with different backgrounds, lighting, and partial occlusions.

\subsection*{The leakage problem}

A critical but often underappreciated methodological issue in animal re-identification
research is \emph{identity leakage}: when images are split randomly into train/val/test
at the image level rather than the identity level, the same individual can appear in all
three partitions.  The model therefore sees every identity during training and the test
protocol measures closed-set retrieval rather than genuine open-set generalisation.
We demonstrate that a Siamese network trained and evaluated under the random-pair protocol
achieves $F_1 = 0.99$, whereas the same network evaluated under a strict identity-disjoint
protocol on the same images yields $F_1 = 0.71$.  The lower number is the scientifically
correct one.

\subsection*{Contributions}

This work makes the following contributions:
\begin{enumerate}[leftmargin=1.6em, itemsep=2pt]
  \item We introduce and validate an \textbf{identity-disjoint evaluation framework} for
        bat face recognition, enforced at the manifest level with cryptographic integrity
        checking, and demonstrate that it is the correct protocol for open-set individual
        recognition.
  \item We conduct a \textbf{systematic multi-factor comparison} of pair-based and
        embedding-based deep recognition architectures across two bat species, three
        background conditions, and controlled identity counts, using real video data from
        tracking studies.
  \item We show that, at the scale of ecological datasets ($\le 13$ training identities),
        the pair-based Siamese network is the only family that verifies significantly
        across all conditions, while corrected angular-margin embeddings (ArcFace/AdaFace)
        match or surpass it on coherent backgrounds---characterising the embedding-model
        shortfall as a \textbf{background-sensitivity} effect (collapse on random-composited
        backgrounds) rather than a fixed small-identity ceiling.
  \item We release a \textbf{reproducible 13-package pipeline} (\batcli) that automates
        the full workflow from raw video to PDF experimental report, with Hydra
        configuration, MLflow tracking, and Optuna hyperparameter search.
\end{enumerate}

The remainder of this paper is organised as follows.
Section~\ref{sec:related} reviews related work.
Section~\ref{sec:dataset} describes the datasets and the identity-disjoint splitter.
Section~\ref{sec:methods} details the three model families.
Section~\ref{sec:evaluation} formalises the evaluation protocol and metrics.
Section~\ref{sec:experiments} presents experimental results.
Section~\ref{sec:discussion} analyses the findings and limitations.
Section~\ref{sec:conclusion} concludes.
